\documentclass{article}

% if you need to pass options to natbib, use, e.g.:
%     \PassOptionsToPackage{numbers, compress}{natbib}
% before loading neurips_2026

% The authors should use one of these tracks.
% Before accepting by the NeurIPS conference, select one of the options below.
% 0. "default" for submission
\usepackage{neurips_2026}
% the "default" option is equal to the "main" option, which is used for the Main Track with double-blind reviewing.
% 1. "main" option is used for the Main Track
%  \usepackage[main]{neurips_2026}
% 2. "position" option is used for the Position Paper Track
%  \usepackage[position]{neurips_2026}
% 3. "eandd" option is used for the Evaluations & Datasets Track
 % \usepackage[eandd]{neurips_2026}
 % if you want to opt-in for a de-anomymized submission (not recommended, but OK):
 %\usepackage[eandd, nonanonymous]{neurips_2026}
% 4. "creativeai" option is used for the Creative AI Track
%  \usepackage[creativeai]{neurips_2026}
% 5. "sglblindworkshop" option is used for the Workshop with single-blind reviewing
 % \usepackage[sglblindworkshop]{neurips_2026}
% 6. "dblblindworkshop" option is used for the Workshop with double-blind reviewing
%  \usepackage[dblblindworkshop]{neurips_2026}

% After being accepted, the authors should add "final" behind the track to compile a camera-ready version.
% 1. Main Track
 % \usepackage[main, final]{neurips_2026}
% 2. Position Paper Track
%  \usepackage[position, final]{neurips_2026}
% 3. Evaluations & Datasets Track
 % \usepackage[eandd, final]{neurips_2026}
% 4. Creative AI Track
%  \usepackage[creativeai, final]{neurips_2026}
% 5. Workshop with single-blind reviewing
%  \usepackage[sglblindworkshop, final]{neurips_2026}
% 6. Workshop with double-blind reviewing
%  \usepackage[dblblindworkshop, final]{neurips_2026}
% Note. For the workshop paper template, both \title{} and \workshoptitle{} are required, with the former indicating the paper title shown in the title and the latter indicating the workshop title displayed in the footnote.
% For workshops (5., 6.), the authors should add the name of the workshop, "\workshoptitle" command is used to set the workshop title.
% \workshoptitle{WORKSHOP TITLE}

% "preprint" option is used for arXiv or other preprint submissions
 % \usepackage[preprint]{neurips_2026}

% to avoid loading the natbib package, add option nonatbib:
%    \usepackage[nonatbib]{neurips_2026}

\usepackage[utf8]{inputenc} % allow utf-8 input
\usepackage[T1]{fontenc}    % use 8-bit T1 fonts
\usepackage{hyperref}       % hyperlinks
\usepackage{url}            % simple URL typesetting
\usepackage{booktabs}       % professional-quality tables
\usepackage{amsfonts}       % blackboard math symbols
\usepackage{nicefrac}       % compact symbols for 1/2, etc.
\usepackage{microtype}      % microtypography
\usepackage{xcolor}         % colors

% Note. For the workshop paper template, both \title{} and \workshoptitle{} are required, with the former indicating the paper title shown in the title and the latter indicating the workshop title displayed in the footnote. 
\title{Formatting Instructions For NeurIPS 2026}


% The \author macro works with any number of authors. There are two commands
% used to separate the names and addresses of multiple authors: \And and \AND.
%
% Using \And between authors leaves it to LaTeX to determine where to break the
% lines. Using \AND forces a line break at that point. So, if LaTeX puts 3 of 4
% authors names on the first line, and the last on the second line, try using
% \AND instead of \And before the third author name.


\author{%
  David S.~Hippocampus\thanks{Use footnote for providing further information
    about author (webpage, alternative address)---\emph{not} for acknowledging
    funding agencies.} \\
  Department of Computer Science\\
  Cranberry-Lemon University\\
  Pittsburgh, PA 15213 \\
  \texttt{hippo@cs.cranberry-lemon.edu} \\
  % examples of more authors
  % \And
  % Coauthor \\
  % Affiliation \\
  % Address \\
  % \texttt{email} \\
  % \AND
  % Coauthor \\
  % Affiliation \\
  % Address \\
  % \texttt{email} \\
  % \And
  % Coauthor \\
  % Affiliation \\
  % Address \\
  % \texttt{email} \\
  % \And
  % Coauthor \\
  % Affiliation \\
  % Address \\
  % \texttt{email} \\
}


\begin{document}


\maketitle


\begin{abstract}
  The abstract paragraph should be indented \nicefrac{1}{2}~inch (3~picas) on both the left- and right-hand margins. Use 10~point type, with a vertical spacing (leading) of 11~points. The word \textbf{Abstract} must be centered, bold, and in point size 12. Two line spaces precede the abstract. The abstract must be limited to one paragraph.
\end{abstract}



\section{Introduction}

Word-of-mouth and viral marketing techniques have become dominant strategies for product promotion in social networks. By selecting a small set of influential seed users, companies can trigger cascades of adoption along social connections, a problem commonly referred to as influence maximization (IM). Over the past two decades, IM has been extensively studied under classical models such as the Independent Cascade (IC) and Linear Threshold (LT) models, where activated nodes attempt to influence their neighbors in a broadcast manner. However, many real-world campaigns deviate from this paradigm. For example, when a retailer distributes digital coupons, a user who chooses not to redeem or discard the coupon may forward it to at most one other user. A similar pattern arises in invitation-code campaigns, where an inactive user can pass the code to only one specific friend. In both scenarios, the promotional content propagates via a sequential, one-to-one diffusion process. We refer to this setting as sequential influence diffusion (SID), which admits a formulation closely related to random walks on the social graph. Despite its practical relevance, this setting remains underexplored, with existing IM frameworks inherently tailored to broadcast diffusion and unable to capture such sequential propagation mechanisms.

This sequential propagation fundamentally alters both the diffusion mechanism and the objective landscape. First, the diffusion model is structurally different. In standard IC and LT models, an activated node attempts to influence all its neighbors simultaneously, resulting in a branching cascade that fans out rapidly across the network. In SID, however, each piece of promotional content is passed to at most one neighbor, tracing a chain-like propagation process through the network. Second, the definition of activation is strictly decoupled from mere exposure. In classical IM, a node is considered activated as soon as the message reaches it. In coupon or invitation-code scenarios, however, a node that simply forwards the coupon without using it is not considered activated; only a node that ultimately redeems the coupon counts as an activated user. Consequently, the intermediate nodes along the forwarding chain contribute nothing to the final activation. Third, since multiple coupons propagate independently, the same user may redeem more than one coupon, yet is counted as activated only once. The primary goal is therefore to maximize the number of distinct users who ultimately redeem at least one coupon, as each newly activated user represents a potential recurring customer. Beyond this, when multiple seeding strategies yield the same activation coverage, the total number of coupons redeemed also serves as an additional metric. These differences collectively fall outside the modeling scope of existing IM frameworks and call for a new problem formulation.

To address these challenges, we formalize the \emph{Sequential Influence Diffusion Influence Maximization (SIDIM)} problem. In our model, $k$ coupons are distributed to $k$ seed nodes, and each coupon propagates independently through the network as a random walk. At each node along the walk, one of three outcomes occurs: the node redeems the coupon and becomes activated, the node discards the coupon, or the node forwards the coupon to exactly one randomly selected neighbor. In the first two cases, the random walk terminates; in the third, it continues. The objective is to select $k$ seed nodes that maximize the expected number of distinct users who redeem at least one coupon. Since the $k$ coupons propagate independently, each coupon induces its own possible world along the nodes it may visit. The overall diffusion outcome is then determined by the joint configuration of $k$ independent possible worlds. We first prove that SIDIM is NP-hard, confirming the inherent computational intractability of the problem. We then analyze the structural properties of the objective function through the possible-world framework and establish that the expected adoption spread is monotone and submodular. This key property enables the application of a greedy algorithm with a $(1-1/e)$-approximation guarantee. To make the greedy algorithm scalable, we develop a novel Reverse Influence Sampling (RIS) framework tailored to the sequential diffusion setting. Unlike classical RR set construction, where a single RR set captures all potential influence paths to a target node, our framework jointly samples $k$ reverse reachable sets for each target --- one per coupon, corresponding to one per possible world --- reflecting the independent random walks in the diffusion process. Building on this RIS framework, we design a near-linear time seed selection algorithm with a $(1-1/e-\epsilon)$-approximation guarantee.

Our contributions are summarized as follows. First, we initiate the formal study of the SIDIM problem, a practically motivated yet theoretically underexplored setting that captures the single-path, multi-coupon nature of real-world sequential promotional campaigns. Second, we prove that SIDIM is NP-hard and establish that the adoption spread function is monotone and submodular, admitting a greedy algorithm with a $(1-1/e)$-approximation guarantee. Third, we develop a tailored RIS scheme that jointly generates $k$ reverse reachable sets per root node with provable unbiasedness, and build upon it a near-linear time algorithm with a $(1-1/e-\epsilon)$-approximation guarantee. Fourth, extensive experiments on real-world social networks demonstrate that our method significantly outperforms baselines adapted from classical IM, validating both the importance of the sequential diffusion formulation and the effectiveness of our algorithmic framework.


\section{Submission of papers to NeurIPS 2026}


Please read the instructions below carefully and follow them faithfully.


\subsection{Style}


Papers to be submitted to NeurIPS 2026 must be prepared according to the
instructions presented here. Papers may only be up to {\bf nine} pages long, including figures. \textbf{Papers that exceed the page limit will not be reviewed (or in any other way considered) for presentation at the conference.}
Additional pages \emph{containing acknowledgments, references, checklist, and optional technical appendices} do not count as content pages. 


The margins in 2026 are the same as those in previous years.


Authors are required to use the NeurIPS \LaTeX{} style files obtainable at the NeurIPS website as indicated below. Please make sure you use the current files and not previous versions. Tweaking the style files may be grounds for desk rejection.


\subsection{Retrieval of style files}


The style files for NeurIPS and other conference information are available on the website at
\begin{center}
  \url{https://neurips.cc}.
\end{center}
% The file \verb+neurips_2026.pdf+ contains these instructions and illustrates the various formatting requirements your NeurIPS paper must satisfy.


The only supported style file for NeurIPS 2026 is \verb+neurips_2026.sty+, rewritten for \LaTeXe{}. \textbf{Previous style files for \LaTeX{} 2.09,  Microsoft Word, and RTF are no longer supported.}


The \LaTeX{} style file contains three optional arguments: 
\begin{itemize}
\item \verb+final+, which creates a camera-ready copy, 
\item \verb+preprint+, which creates a preprint for submission to, e.g., arXiv, \item \verb+nonatbib+, which will not load the \verb+natbib+ package for you in case of package clash.
\end{itemize}


\paragraph{Preprint option}
If you wish to post a preprint of your work online, e.g., on arXiv, using the NeurIPS style, please use the \verb+preprint+ option. This will create a nonanonymized version of your work with the text ``Preprint. Work in progress.'' in the footer. This version may be distributed as you see fit, as long as you do not say which conference it was submitted to. Please \textbf{do not} use the \verb+final+ option, which should \textbf{only} be used for papers accepted to NeurIPS.


At submission time, please omit the \verb+final+ and \verb+preprint+ options. This will anonymize your submission and add line numbers to aid review. Please do \emph{not} refer to these line numbers in your paper as they will be removed during generation of camera-ready copies.


The file \verb+neurips_2026.tex+ may be used as a ``shell'' for writing your paper. All you have to do is replace the author, title, abstract, and text of the paper with your own.


The formatting instructions contained in these style files are summarized in Sections \ref{gen_inst}, \ref{headings}, and \ref{others} below.
















\section{Problem Definition}

Let $G = (V, E)$ be a directed graph with $n = |V|$ nodes and $m = |E|$ edges, where a node $u \in V$ represents a user and a directed edge $(u, v) \in E$ indicates that $u$ may forward a coupon to $v$. Throughout this paper, we assume $G$ is strongly connected. Given a seed set $S = \{s_1, s_2, \ldots, s_k\} \subseteq V$ of size $k$, a retailer distributes $k$ coupons, one to each seed user, and each coupon then propagates independently through the network according to the diffusion process described below.

Consider first the propagation of a single coupon. When a coupon arrives at a node $u$, one of three mutually exclusive outcomes occurs. With probability $p_u^a$, user $u$ \emph{adopts} the coupon: the coupon is redeemed, $u$ is marked as activated, and the propagation terminates at $u$. With probability $p_u^d$, user $u$ \emph{discards} the coupon: the coupon is removed from the network without activating anyone. With the remaining probability $p_u^t = 1 - p_u^a - p_u^d$, user $u$ \emph{transfers} the coupon to exactly one out-neighbor $v$, selected according to an edge-level forwarding parameter $p(u, v)$ satisfying $\sum_{v: (u, v) \in E} p(u, v) = p_u^t$. A crucial feature of this process, which sharply distinguishes our setting from classical IM models, is that the coupon itself is \emph{non-replicable}: unlike a piece of information in the IC or LT model that can be copied and broadcast to all neighbors simultaneously, a coupon exists as a single instance in the network and moves to at most one location at each step.

Building on this single-coupon process, the full diffusion proceeds by running $k$ such processes in parallel, one for each coupon-seed pair, with all stochastic decisions drawn independently across coupons. Because the $k$ coupons evolve independently, a node $u$ may be visited by multiple coupons during the overall diffusion, and may even adopt more than one of them. Yet from the retailer's standpoint, the same user being reached by several coupons is no more valuable than being reached by one. Accordingly, a user $u$ is considered ultimately \emph{activated} if at least one of the $k$ coupons is adopted at $u$, and is counted exactly once in the final outcome regardless of how many coupons reach or are redeemed by $u$.

\textit{Remark.} It is worth highlighting a subtle modeling distinction from classical IM. In the IC model, the activation along each edge is governed by an independent Bernoulli trial, and a node may be influenced through multiple edges in parallel. In our model, by contrast, the three outcomes a node faces upon receiving a coupon are \emph{mutually exclusive and competing for the same decision moment}: the user must choose between redeeming, discarding, or forwarding to a particular neighbor, and these choices share a common probability budget summing to one. This coupling between the three outcomes, together with the non-replicable nature of the coupon, is what gives rise to the chain-like propagation behavior characteristic of our setting.

Having specified the diffusion process, we now formalize the optimization objective. Let $\mu(u, S) \in \{0, 1\}$ denote whether user $u$ is activated under the seed set $S$, and let $I_C(S) = \mathbb{E}[\sum_{u \in V} \mu(u, S)]$ denote the expected number of activated users, taken over the randomness of the propagation process $C$. We refer to $I_C(S)$ as the \emph{adoption spread} of $S$. The problem in this paper is formalized as:

\textbf{Problem Definition 1 (Coupon Influence Maximization)}
Given a directed graph $G = (V, E)$ with node parameters $\{(p_u^a, p_u^d)\}_{u \in V}$, edge forwarding parameters $\{p(u, v)\}_{(u, v) \in E}$, and an integer $k$, find a size-$k$ seed set $S^*$ that maximizes the expected adoption spread, i.e.,
\begin{equation}
S^* = \arg\max_{S \subseteq V, \, |S| = k} I_C(S).
\end{equation}











\section{Complexity Analysis}

In this section, we analyze the computational complexity and structural properties of the Coupon Influence Maximization problem, which together motivate our algorithmic design in the next section.

\begin{theorem}
\label{thm:nphard}
The Coupon Influence Maximization problem is NP-hard.
\end{theorem}

\begin{proof}[Proof sketch]
The classical IM problem under the IC model can be cast as a special instance of our problem by setting $p_u^d = 0$ for every node $u$ and choosing the forwarding parameters $\{p(u, v)\}$ appropriately so that the resulting single-coupon walk reproduces an IC-style cascade restricted to a single trajectory. Since classical IM is NP-hard~\cite{kempe2003maximizing}, our problem is at least as hard. The full reduction is provided in the appendix.
\end{proof}

To analyze the structural properties of the adoption spread $I_C(S)$, it is convenient to adopt a \emph{possible-world} perspective on the diffusion. For each coupon $j \in \{1, \ldots, k\}$ and each node $u \in V$, the stochastic outcome at $u$ when coupon $j$ visits it (adoption, discarding, or transfer to a specific neighbor) can be encoded by a single random variable, jointly determining $u$'s behavior under coupon $j$. The collection of these random variables across all nodes constitutes a possible world $W_j$ for coupon $j$, and the joint configuration $\mathcal{W} = (W_1, \ldots, W_k)$ of $k$ independent possible worlds fully determines the diffusion outcome. Within a fixed $\mathcal{W}$, the activation indicator $\mu(u, S)$ becomes a deterministic function of $S$, and the adoption spread is recovered by taking expectation: $I_C(S) = \mathbb{E}_{\mathcal{W}}[\sum_u \mu(u, S)]$.

\begin{theorem}
\label{thm:monosub}
The adoption spread $I_C(S)$ is monotone non-decreasing and submodular in $S$.
\end{theorem}

\begin{proof}[Proof sketch]
Fix an arbitrary realization $\mathcal{W}$ of the $k$ possible worlds. Under this realization, a user $u$ is activated by seed set $S$ if and only if at least one seed $s_j \in S$ can reach $u$ via a deterministic path of transfers determined by $\mathcal{W}$ and ultimately triggers an adoption at $u$. Equivalently, for each user $u$ there is a deterministic set $T_u(\mathcal{W}) \subseteq V$ of nodes that, when chosen as a seed, would activate $u$ under $\mathcal{W}$, and $\mu(u, S) = \mathbb{I}[S \cap T_u(\mathcal{W}) \neq \emptyset]$.

Monotonicity is immediate: enlarging $S$ can only increase the chance that $S \cap T_u(\mathcal{W}) \neq \emptyset$. For submodularity, observe that for any fixed $\mathcal{W}$, the function $S \mapsto \mathbb{I}[S \cap T_u(\mathcal{W}) \neq \emptyset]$ is a coverage function and hence submodular in $S$. Summing over $u \in V$ preserves submodularity, and taking the expectation over $\mathcal{W}$ preserves both monotonicity and submodularity, since these properties are closed under non-negative linear combinations. Therefore $I_C(S)$ is monotone and submodular. The detailed proof is provided in the appendix.
\end{proof}

Theorems~\ref{thm:nphard} and~\ref{thm:monosub} together characterize the Coupon Influence Maximization problem as a monotone submodular maximization problem subject to a cardinality constraint. By the classical result of Nemhauser et al.~\cite{nemhauser1978analysis}, a simple greedy algorithm that iteratively adds the node with the largest marginal gain achieves a $(1 - 1/e)$-approximation of the optimum. However, evaluating the marginal gain exactly is intractable, since computing $I_C(S)$ for a given $S$ already requires reasoning over all possible worlds. In the next section, we develop a Reverse Influence Sampling framework tailored to the coupon diffusion process, which estimates marginal gains efficiently and yields a near-linear time algorithm with a $(1 - 1/e - \epsilon)$-approximation guarantee.








































\section{General formatting instructions}
\label{gen_inst}


The text must be confined within a rectangle 5.5~inches (33~picas) wide and
9~inches (54~picas) long. The left margin is 1.5~inch (9~picas).  Use 10~point
type with a vertical spacing (leading) of 11~points.  Times New Roman is the
preferred typeface throughout, and will be selected for you by default.
Paragraphs are separated by \nicefrac{1}{2}~line space (5.5 points), with no
indentation.


The paper title should be 17~point, initial caps/lower case, bold, centered
between two horizontal rules. The top rule should be 4~points thick and the
bottom rule should be 1~point thick. Allow \nicefrac{1}{4}~inch space above and
below the title to rules. All pages should start at 1~inch (6~picas) from the
top of the page.


For the final version, authors' names are set in boldface, and each name is
centered above the corresponding address. The lead author's name is to be listed
first (left-most), and the co-authors' names (if different address) are set to
follow. If there is only one co-author, list both author and co-author side by
side.


Please pay special attention to the instructions in Section \ref{others}
regarding figures, tables, acknowledgments, and references.

\section{Headings: first level}
\label{headings}


All headings should be lower case (except for first word and proper nouns),
flush left, and bold.


First-level headings should be in 12-point type.


\subsection{Headings: second level}


Second-level headings should be in 10-point type.


\subsubsection{Headings: third level}


Third-level headings should be in 10-point type.


\paragraph{Paragraphs}


There is also a \verb+\paragraph+ command available, which sets the heading in
bold, flush left, and inline with the text, with the heading followed by 1\,em
of space.


\section{Citations, figures, tables, references}
\label{others}


These instructions apply to everyone.


\subsection{Citations within the text}


The \verb+natbib+ package will be loaded for you by default.  Citations may be
author/year or numeric, as long as you maintain internal consistency.  As to the
format of the references themselves, any style is acceptable as long as it is
used consistently.


The documentation for \verb+natbib+ may be found at
\begin{center}
  \url{http://mirrors.ctan.org/macros/latex/contrib/natbib/natnotes.pdf}
\end{center}
Of note is the command \verb+\citet+, which produces citations appropriate for
use in inline text.  For example,
\begin{verbatim}
   \citet{hasselmo} investigated\dots
\end{verbatim}
produces
\begin{quote}
  Hasselmo, et al.\ (1995) investigated\dots
\end{quote}


If you wish to load the \verb+natbib+ package with options, you may add the
following before loading the \verb+neurips_2026+ package:
\begin{verbatim}
   \PassOptionsToPackage{options}{natbib}
\end{verbatim}


If \verb+natbib+ clashes with another package you load, you can add the optional
argument \verb+nonatbib+ when loading the style file:
\begin{verbatim}
   \usepackage[nonatbib]{neurips_2026}
\end{verbatim}


As submission is double blind, refer to your own published work in the third
person. That is, use ``In the previous work of Jones et al.\ [4],'' not ``In our
previous work [4].'' If you cite your other papers that are not widely available
(e.g., a journal paper under review), use anonymous author names in the
citation, e.g., an author of the form ``A.\ Anonymous'' and include a copy of the anonymized paper in the supplementary material.


\subsection{Footnotes}


Footnotes should be used sparingly.  If you do require a footnote, indicate
footnotes with a number\footnote{Sample of the first footnote.} in the
text. Place the footnotes at the bottom of the page on which they appear.
Precede the footnote with a horizontal rule of 2~inches (12~picas).


Note that footnotes are properly typeset \emph{after} punctuation
marks.\footnote{As in this example.}


\subsection{Figures}


\begin{figure}
  \centering
  \fbox{\rule[-.5cm]{0cm}{4cm} \rule[-.5cm]{4cm}{0cm}}
  \caption{Sample figure caption. Explain what the figure shows and add a key take-away message to the caption.}
\end{figure}


All artwork must be neat, clean, and legible. Lines should be dark enough for
 reproduction purposes. The figure number and caption always appear after the
figure. Place one line space before the figure caption and one line space after
the figure. The figure caption should be lower case (except for the first word and proper nouns); figures are numbered consecutively.


You may use color figures.  However, it is best for the figure captions and the
paper body to be legible if the paper is printed in either black/white or in
color.


\subsection{Tables}


All tables must be centered, neat, clean, and legible.  The table number and
title always appear before the table.  See Table~\ref{sample-table}.


Place one line space before the table title, one line space after the
table title, and one line space after the table. The table title must
be lower case (except for the first word and proper nouns); tables are
numbered consecutively.


Note that publication-quality tables \emph{do not contain vertical rules}. We
strongly suggest the use of the \verb+booktabs+ package, which allows for
typesetting high-quality, professional tables:
\begin{center}
  \url{https://www.ctan.org/pkg/booktabs}
\end{center}
This package was used to typeset Table~\ref{sample-table}.


\begin{table}
  \caption{Sample table caption. Explain what the table shows and add a key take-away message to the caption.}
  \label{sample-table}
  \centering
  \begin{tabular}{lll}
    \toprule
    \multicolumn{2}{c}{Part}                   \\
    \cmidrule(r){1-2}
    Name     & Description     & Size ($\mu$m) \\
    \midrule
    Dendrite & Input terminal  & $\approx$100     \\
    Axon     & Output terminal & $\approx$10      \\
    Soma     & Cell body       & up to $10^6$  \\
    \bottomrule
  \end{tabular}
\end{table}

\subsection{Math}
Note that display math in bare TeX commands will not create correct line numbers for submission. Please use LaTeX (or AMSTeX) commands for unnumbered display math. (You really shouldn't be using \$\$ anyway; see \url{https://tex.stackexchange.com/questions/503/why-is-preferable-to} and \url{https://tex.stackexchange.com/questions/40492/what-are-the-differences-between-align-equation-and-displaymath} for more information.)

\subsection{Final instructions}

Do not change any aspects of the formatting parameters in the style files.  In
particular, do not modify the width or length of the rectangle the text should
fit into, and do not change font sizes. Please note that pages should be
numbered.


\section{Preparing PDF files}


Please prepare submission files with paper size ``US Letter,'' and not, for
example, ``A4.''


Fonts were the main cause of problems in the past years. Your PDF file must only
contain Type 1 or Embedded TrueType fonts. Here are a few instructions to
achieve this.


\begin{itemize}


\item You should directly generate PDF files using \verb+pdflatex+.


\item You can check which fonts a PDF files uses.  In Acrobat Reader, select the
  menu Files$>$Document Properties$>$Fonts and select Show All Fonts. You can
  also use the program \verb+pdffonts+ which comes with \verb+xpdf+ and is
  available out-of-the-box on most Linux machines.


\item \verb+xfig+ ``patterned'' shapes are implemented with bitmap fonts.  Use
  "solid" shapes instead.


\item The \verb+\bbold+ package almost always uses bitmap fonts.  You should use
  the equivalent AMS Fonts:
\begin{verbatim}
   \usepackage{amsfonts}
\end{verbatim}
followed by, e.g., \verb+\mathbb{R}+, \verb+\mathbb{N}+, or \verb+\mathbb{C}+
for $\mathbb{R}$, $\mathbb{N}$ or $\mathbb{C}$.  You can also use the following
workaround for reals, natural and complex:
\begin{verbatim}
   \newcommand{\RR}{I\!\!R} %real numbers
   \newcommand{\Nat}{I\!\!N} %natural numbers
   \newcommand{\CC}{I\!\!\!\!C} %complex numbers
\end{verbatim}
Note that \verb+amsfonts+ is automatically loaded by the \verb+amssymb+ package.


\end{itemize}


If your file contains type 3 fonts or non embedded TrueType fonts, we will ask
you to fix it.


\subsection{Margins in \LaTeX{}}


Most of the margin problems come from figures positioned by hand using
\verb+\special+ or other commands. We suggest using the command
\verb+\includegraphics+ from the \verb+graphicx+ package. Always specify the
figure width as a multiple of the line width as in the example below:
\begin{verbatim}
   \usepackage[pdftex]{graphicx} ...
   \includegraphics[width=0.8\linewidth]{myfile.pdf}
\end{verbatim}
See Section 4.4 in the graphics bundle documentation
(\url{http://mirrors.ctan.org/macros/latex/required/graphics/grfguide.pdf})


A number of width problems arise when \LaTeX{} cannot properly hyphenate a
line. Please give LaTeX hyphenation hints using the \verb+\-+ command when
necessary.

\begin{ack}
Use unnumbered first level headings for the acknowledgments. All acknowledgments
go at the end of the paper before the list of references. Moreover, you are required to declare
funding (financial activities supporting the submitted work) and competing interests (related financial activities outside the submitted work).
More information about this disclosure can be found at: \url{https://neurips.cc/Conferences/2026/PaperInformation/FundingDisclosure}.


Do {\bf not} include this section in the anonymized submission, only in the final paper. You can use the \texttt{ack} environment provided in the style file to automatically hide this section in the anonymized submission.
\end{ack}

\section*{References}


References follow the acknowledgments in the camera-ready paper. Use unnumbered first-level heading for
the references. Any choice of citation style is acceptable as long as you are
consistent. It is permissible to reduce the font size to \verb+small+ (9 point)
when listing the references.
Note that the Reference section does not count towards the page limit.
\medskip


{
\small


[1] Alexander, J.A.\ \& Mozer, M.C.\ (1995) Template-based algorithms for
connectionist rule extraction. In G.\ Tesauro, D.S.\ Touretzky and T.K.\ Leen
(eds.), {\it Advances in Neural Information Processing Systems 7},
pp.\ 609--616. Cambridge, MA: MIT Press.


[2] Bower, J.M.\ \& Beeman, D.\ (1995) {\it The Book of GENESIS: Exploring
  Realistic Neural Models with the GEneral NEural SImulation System.}  New York:
TELOS/Springer--Verlag.


[3] Hasselmo, M.E., Schnell, E.\ \& Barkai, E.\ (1995) Dynamics of learning and
recall at excitatory recurrent synapses and cholinergic modulation in rat
hippocampal region CA3. {\it Journal of Neuroscience} {\bf 15}(7):5249-5262.
}


%%%%%%%%%%%%%%%%%%%%%%%%%%%%%%%%%%%%%%%%%%%%%%%%%%%%%%%%%%%%

\appendix

\section{Technical appendices and supplementary material}
Technical appendices with additional results, figures, graphs, and proofs may be submitted with the paper submission before the full submission deadline (see above). You can upload a ZIP file for videos or code, but do not upload a separate PDF file for the appendix. There is no page limit for the technical appendices. 

Note: Think of the appendix as ``optional reading'' for reviewers. The paper must be able to stand alone without the appendix; for example, adding critical experiments that support the main claims to an appendix is inappropriate. 

%%%%%%%%%%%%%%%%%%%%%%%%%%%%%%%%%%%%%%%%%%%%%%%%%%%%%%%%%%%%

\newpage
\input{checklist.tex}


\end{document}